\documentclass[tikz]{standalone}
\usepackage{../sty/phylogenetic-diagrams} % Include the style file containing the code defining the structure and style of the document

\begin{document}
\begin{tikzpicture}[node distance=2.5cm]
	\node (time_origin) at (-4,0) {$t_{\Omega}$};
	\node (time_archetype) at (-4,-2) {$t_{\omega}$};
	\node (time_alpha) at (-4,-3) {$t_{\alpha}$};
	\node (time_beta) at (-4,-3.5) {$t_{\beta}$};
	\node (time_gamma) at (-4,-5.5) {$t_{\gamma}$};
	\node (time_a) at (-4,-4.5) {$t_{\mathrm{a}}$};
	\node (time_b1) at (-4,-6) {$t_{\mathrm{b}_1}$};
	\node (time_b2) at (-4,-6.5) {$t_{\mathrm{b}_2}$};
	\node (time_c1) at (-4,-7) {$t_{\mathrm{c}_1}$};
	\node (time_c2) at (-4,-8) {$t_{\mathrm{c}_2}$};
	\node (time_0) at (-4,-9) {$0$};
	\draw[dashed] (time_origin) -- (time_archetype) -- (time_alpha) -- (time_beta) -- (time_a) -- (time_gamma) -- (time_b1) -- (time_b2) -- (time_c1) -- (time_c2) -- (time_0);
	\node[hypotheticalWit] (origin) at (0,0) {$\Omega$};
	\node[hypotheticalWit] (archetype) at (0,-2) {$\omega$};
	\node[hypotheticalWit] (alpha) at (-2,-3) {$\alpha$};
	\node[hypotheticalWit] (beta) at (2,-3.5) {$\beta$};
	\node[extantWit] (a) at (-3,-4.5) {$\mathrm{a}$};
	\node[hypotheticalWit] (gamma) at (-1,-5.5) {$\gamma$};
	\node[extantWit] (b1) at (1,-6) {$\mathrm{b}_1$};
	\node[extantWit] (b2) at (3,-6.5) {$\mathrm{b}_2$};
	\node[extantWit] (c1) at (-1.5,-7) {$\mathrm{c}_1$};
	\node[extantWit] (c2) at (-0.5,-8) {$\mathrm{c}_2$};
	\draw[undirected] (origin.south) -| (archetype.north);
	\draw[undirected] (archetype.west) -| (alpha.north);
	\draw[undirected] (archetype.east) -| (beta.north);
	\draw[undirected] (alpha.west) -| (a.north);
	\draw[undirected] (alpha.east) -| (gamma.north);
	\draw[undirected] (beta.west) -| (b1.north);
	\draw[undirected] (beta.east) -| (b2.north);
	\draw[undirected] (gamma.west) -| (c1.north);
	\draw[undirected] (gamma.east) -| (c2.north);
\end{tikzpicture}
\end{document}